\begin{table*}[t]
\small
\centering
\begin{tabular}{@{}p{.20\textwidth}p{.73\textwidth}@{}}
\textbf{Reporting item} & \textbf{Information to provide} \\
\toprule
Target population & Who or what is the study about? Describe observed users, communities, languages, dates, and exclusions; distinguish the sample from the target population. \\
\rowcollight Question and estimand & State the descriptive quantity or causal contrast, unit of analysis, exposure, outcome, and time horizon. Specify whose effect is estimated. \\
Constructs and measures & Explain why each trace or label measures the intended construct. Report validation, disagreements, missingness, and relevant subgroup or period differences. \\
\rowcollight Collection and provenance & Record interfaces, sampling rules, access restrictions, deletion or attrition, and data and model versions. Document content-production processes where known. \\
Assignment and timing & Describe how exposure arises, the comparison group, and the ordering of covariates, treatment, and outcomes. Identify concurrent events and possible spillovers. \\
\rowcollight Adjustment and overlap & Justify pretreatment covariates, including text features. Report balance, overlap, exclusions, and any change in the target population after adjustment. \\
Sensitivity and uncertainty & State assumptions and plausible failures. Use appropriate sensitivity analyses, alternative measures, placebo or negative-control comparisons, and uncertainty estimates. \\
\rowcollight Interpretation and reuse & Separate prediction, association, and causal interpretation. State limits of transfer and provide reusable code, definitions, and access information when possible. \\
\bottomrule
\end{tabular}
\caption{A reporting checklist for studies using social media observations. Items should be applied to the study's question and design, with non-applicable items explained. The checklist is a reporting aid, not a quality score or a substitute for validation.}
\label{tab:checklist}
\end{table*}
